\section{Integer witnesses in every rank at least five}\label{sec:witness}
\subsection{Six rank-one matrices and a five-term tensor}

Identify $\C^3\otimes\C^3$ with $3\times3$ matrices, so that
$a\otimes b$ corresponds to $ab^{\mathsf T}$. Set
\begin{equation}\label{eq:M}
 M=\begin{pmatrix}1&1&1\\1&2&3\\1&4&9\end{pmatrix}.
\end{equation}
Let $r_i$ be its $i$th row, $m_j$ its $j$th column, and
$e_1,e_2,e_3$ the coordinate vectors. Define
\begin{equation}\label{eq:Q}
 Q_i=e_i r_i\quad(1\leq i\leq3),\qquad
 Q_{3+j}=-m_j e_j^{\mathsf T}\quad(1\leq j\leq3).
\end{equation}
Thus $Q_s=a_s b_s^{\mathsf T}$ for the following integer factors:
\begin{center}
\begin{tabular}{c@{\qquad}c@{\qquad}c}
\toprule
$s$&$a_s^{\mathsf T}$&$b_s^{\mathsf T}$\\
\midrule
$1$&$(1,0,0)$&$(1,1,1)$\\
$2$&$(0,1,0)$&$(1,2,3)$\\
$3$&$(0,0,1)$&$(1,4,9)$\\
$4$&$(-1,-1,-1)$&$(1,0,0)$\\
$5$&$(-1,-2,-4)$&$(0,1,0)$\\
$6$&$(-1,-3,-9)$&$(0,0,1)$\\
\bottomrule
\end{tabular}
\end{center}
They are nonzero and pairwise nonproportional: the row-supported and
column-supported matrices have different supports, and all entries of
$M$ are nonzero.

\begin{lemma}\label{lem:L}\label{lemma:witness-core}
The only linear relation among $Q_1,\ldots,Q_6$, up to a scalar,
is
\begin{equation}\label{eq:Q-relation}
 Q_1+Q_2+Q_3+Q_4+Q_5+Q_6=0.
\end{equation}
Their span $L$ has dimension five, and any five of these matrices form
a basis of $L$. Every $N\in L$ has an expression
\begin{equation}\label{eq:L-coordinates}
 N_{ij}=M_{ij}(x_i-y_j),\qquad 1\leq i,j\leq3.
\end{equation}
\end{lemma}

\begin{proof}
Summing the first three matrices gives $M$, and summing the last
three gives $-M$. If $\sum_s\alpha_s Q_s=0$, its $(i,j)$ entry is
$M_{ij}(\alpha_i-\alpha_{3+j})=0$. Every $M_{ij}$ is nonzero,
so all six coefficients coincide. The unique relation uses every
matrix with a nonzero coefficient, proving independence of any five.
Formula~\eqref{eq:L-coordinates} is the linear combination with
coefficients $x_1,x_2,x_3,y_1,y_2,y_3$ in~\eqref{eq:Q}.
\end{proof}

Let $f_1,\ldots,f_5$ be the standard basis of $\C^5$. Our tensor is
\begin{equation}\label{eq:T-first}
 T=\sum_{i=1}^5 Q_i\otimes f_i
   \in\C^3\otimes\C^3\otimes\C^5.
\end{equation}
This specifies all its integer entries; explicitly, its five matrix
slices are
\begin{equation}\label{eq:slices}
\begin{gathered}
 \begin{pmatrix}1&1&1\\0&0&0\\0&0&0\end{pmatrix},\quad
 \begin{pmatrix}0&0&0\\1&2&3\\0&0&0\end{pmatrix},\quad
 \begin{pmatrix}0&0&0\\0&0&0\\1&4&9\end{pmatrix},\\[3pt]
 \begin{pmatrix}-1&0&0\\-1&0&0\\-1&0&0\end{pmatrix},\quad
 \begin{pmatrix}0&-1&0\\0&-2&0\\0&-4&0\end{pmatrix}.
\end{gathered}
\end{equation}
Relation~\eqref{eq:Q-relation} gives the second decomposition
\begin{equation}\label{eq:T-second}
 T=\sum_{i=1}^4 Q_i\otimes(f_i-f_5)-Q_6\otimes f_5.
\end{equation}
In particular the final summand in~\eqref{eq:T-second} has factors
$(1,3,9)^{\mathsf T}$, $e_3$, and $f_5$; the preceding four have
the factors in the table and the indicated third vectors.

Each expression has five distinct nonzero pure summands. The two
unordered sets are different even if one identifies a pure summand's
different factorizations: the first uses the matrix direction $Q_5$,
whereas the second uses $Q_6$, and these are distinct directions.
In either expression the five third vectors are a basis of $\C^5$,
and the five matrix factors are a basis of $L$.

The third flattening of $T$ has image $L$ and rank five. Every sum
of four pure tensors has all flattening ranks at most four, and the
$5\times5$ minors of a flattening vanish on the Zariski closure of
that set. Consequently the displayed five-term decomposition and
the flattening rank prove
\begin{equation}\label{eq:rank-border-rank}
 \rank(T)=\brr(T)=5.
\end{equation}
Here $\rank(T)$ in~\eqref{eq:rank-border-rank} denotes tensor rank;
matrix ranks are specified as such elsewhere.



\subsection{Six tangent intersections}

Let $\mathcal R$ be the nonzero rank-one matrix cone in
$\C^3\otimes\C^3$. For $Q=ab^{\mathsf T}$ its affine tangent space is
\begin{equation}\label{eq:matrix-tangent}
 T_Q\mathcal R=a(\C^3)^{\mathsf T}+\C^3 b^{\mathsf T}.
\end{equation}
The formula follows by differentiating $ab^{\mathsf T}$; its image
has dimension five, with one-dimensional scaling kernel, and equals
the tangent space of the smooth rank-one locus.

\begin{lemma}\label{lem:tangent-intersections}
For each $s=1,\ldots,6$,
\begin{equation}\label{eq:tangent-intersection}
 L\cap T_{Q_s}\mathcal R=\C Q_s.
\end{equation}
\end{lemma}

\begin{proof}
Fix a row point $Q_i=e_i r_i$, and let $k<\ell$ be the two remaining
row indices. By~\eqref{eq:matrix-tangent}, tangency requires rows
$k$ and $\ell$ of $N$ to be $z_k r_i$ and $z_\ell r_i$.
If $N$ also belongs to $L$, formula~\eqref{eq:L-coordinates} gives
\[
 y_j=x_k-z_k\frac{M_{ij}}{M_{kj}}
     =x_\ell-z_\ell\frac{M_{ij}}{M_{\ell j}}
       \quad(j=1,2,3).
\]
The coefficient vectors in these three equations are independent.
Indeed, define
\[
 D_i=\det\begin{pmatrix}
 1&1&1\\
 M_{i1}/M_{k1}&M_{i2}/M_{k2}&M_{i3}/M_{k3}\\
 M_{i1}/M_{\ell1}&M_{i2}/M_{\ell2}&M_{i3}/M_{\ell3}
 \end{pmatrix}.
\]
Direct rational expansion gives
\begin{equation}\label{eq:row-determinants}
 (D_1,D_2,D_3)=\left(-\frac1{18},\frac13,-2\right).
\end{equation}
Thus $x_k=x_\ell$ and $z_k=z_\ell=0$. All $y_j$ equal $x_k$,
and~\eqref{eq:L-coordinates} leaves only a scalar multiple of $Q_i$.

For a column point $Q_{3+j}$, the minus sign in~\eqref{eq:Q} does
not change the tangent space. Interchange rows and columns in the
preceding argument. With $k<\ell$ the other two column indices,
the required determinant is
\[
 E_j=\det\begin{pmatrix}
 1&1&1\\
 M_{1j}/M_{1k}&M_{2j}/M_{2k}&M_{3j}/M_{3k}\\
 M_{1j}/M_{1\ell}&M_{2j}/M_{2\ell}&M_{3j}/M_{3\ell}
 \end{pmatrix}.
\]
Its values are
\begin{equation}\label{eq:column-determinants}
 (E_1,E_2,E_3)=\left(-\frac1{18},\frac49,-\frac32\right).
\end{equation}
They are again nonzero. The same equations now force every column
other than $j$ to vanish and leave a scalar multiple of $Q_{3+j}$.
This proves~\eqref{eq:tangent-intersection} in all six cases.
\end{proof}


\subsection{From matrix tangencies to tensor immersion}

\begin{lemma}\label{lemma:tangent-criterion}
Let $P_1,\ldots,P_r$ be independent nonzero rank-one matrices in
$A\otimes B$, put $L=\langle P_1,\ldots,P_r\rangle$, and suppose
\[
 L\cap T_{P_i}\mathcal R=\C P_i\quad(1\leq i\leq r),
\]
where $\mathcal R$ is the nonzero rank-one matrix cone.
If $c_1,\ldots,c_r$ are a basis of $C$, then addition on the product
of the $r$ nonzero pure-tensor cones has injective differential at
$(P_i\otimes c_i)_{i=1}^r$.
\end{lemma}

\begin{proof}
Change the basis of $C$ so that $c_i=f_i$. Write a tangent vector at
the $i$th summand as
\[
 \dot P_i\otimes f_i+P_i\otimes\dot c_i,
 \qquad \dot P_i\in T_{P_i}\mathcal R,
 \qquad \dot c_i=\sum_j h_{ji}f_j.
\]
Suppose that the sum of these $r$ tangent vectors is zero. Comparing
coefficients of $f_j$ gives
\[
 \dot P_j+\sum_i h_{ji}P_i=0.
\]
Thus $\dot P_j\in L\cap T_{P_j}\mathcal R=\C P_j$; write
$\dot P_j=\lambda_jP_j$. Independence of the $P_i$ now implies
$h_{ji}=0$ for $i\ne j$ and $h_{jj}=-\lambda_j$. Consequently the
original tangent vector at each individual summand is zero.
The tangent spaces of the pure-tensor cones therefore sum directly.
This argument works intrinsically on the pure-tensor cones, after
removing the factor-coordinate scaling directions.
\end{proof}

Applying the lemma to each decomposition of $T_5:=T$ in
\eqref{eq:T-first} and~\eqref{eq:T-second} gives differential rank
\begin{equation}\label{eq:rank45}
 5(3+3+5-2)=45.
\end{equation}
The sign change from $Q_6$ to $-Q_6$ does not change its tangent space.

\subsection{Two avoidance identities}

Write $A_0=B_0=\C^3$, and denote the span of the six matrices by
$L_0\subset A_0\otimes B_0$. For an integer $t\geq4$ put
\[
 v(t)=(1,t,t^2)^{\mathsf T}.
\]

\begin{lemma}\label{lemma:avoidance}
For every integer $t\geq4$,
\begin{equation}\label{eq:avoidance}
 L_0\cap\bigl(\C v(t)\otimes B_0\bigr)=0,
 \qquad
 L_0\cap\bigl(A_0\otimes\C v(t)\bigr)=0.
\end{equation}
Moreover, the lines $\C v(t)$ are pairwise distinct and avoid all
first- and second-factor lines in the six core matrices.
\end{lemma}

\begin{proof}
Since $M_{ij}=j^{i-1}$, every matrix in $L_0$ has entries
$N_{ij}=j^{i-1}(x_i-y_j)$. First suppose that $N=v(t)b^{\mathsf T}$.
The first row gives $y_j=x_1-b_j$. Substitution into the second row,
with $d=x_1-x_2$, yields
\[
 b_j=\frac{jd}{j-t}\qquad(j=1,2,3).
\]
All denominators are nonzero. Substitution into the third row, with
$e=x_1-x_3$, gives
\[
 (j+t)d=je\qquad(j=1,2,3).
\]
The equations for $j=1,2$ imply $td=0$. Thus $d=e=0$ and $b=0$,
which proves the first intersection statement.

Now suppose $N=a v(t)^{\mathsf T}$. Let
$\mathbf 1=(1,1,1)$ and $y=(y_1,y_2,y_3)$. Row $i$ says
\[
 y=x_i\mathbf1-a_i\frac{v(t)^{\mathsf T}}{r_i},
\]
with entrywise division. The vectors $\mathbf1$, $v(t)^{\mathsf T}$,
and $v(t)^{\mathsf T}/(1,2,3)$ are independent, since
\begin{equation}\label{eq:avoidance-determinant}
 \det\begin{pmatrix}
 1&1&1\\1&t&t^2\\1&t/2&t^2/3
 \end{pmatrix}
 =-\frac{t(t-1)(t-3)}6\ne0.
\end{equation}
The planes obtained from rows one and two therefore meet only in
$\C\mathbf1$. Hence $a_1=a_2=0$ and $y$ is constant.
The vector $(1,t/4,t^2/9)$ from row three is not constant: constancy
would require simultaneously $t=4$ and $t^2=9$. Thus $a_3=0$ as well,
proving the second statement.

Finally, the first coordinate of every $v(t)$ is one, so distinct
integers $t$ give distinct lines. None is a coordinate line.
The other core first-factor lines are those of
$v(1),v(2),v(3)$. The other second-factor vectors are
$(1,1,1)$, $(1,2,3)$ and $(1,4,9)$. Direct comparison with
$(1,t,t^2)$ for $t\geq4$ excludes every such line.
\end{proof}

\subsection{Adding rows and columns}

Fix $r\geq5$ and integers $p,q\geq3$ with $p+q=r+1$. Set
\[
 A'=A_0\oplus E=\C^p,\qquad B'=B_0\oplus F=\C^q,
\]
where $A_0,B_0$ occupy the first three coordinates. Let
$u_1,\ldots,u_{p-3}$ and $w_1,\ldots,w_{q-3}$ be the coordinate bases
of $E$ and $F$. Define the following integer rank-one matrices:
\begin{equation}\label{eq:added-matrices}
 R_s=u_s\otimes v(s+3)\quad(1\leq s\leq p-3),\qquad
 S_t=v(t+3)\otimes w_t\quad(1\leq t\leq q-3).
\end{equation}
Empty lists are permitted. Their span with the core is the direct sum
\begin{equation}\label{eq:extended-L}
 L=L_0\oplus\langle R_1,\ldots,R_{p-3}\rangle
       \oplus\langle S_1,\ldots,S_{q-3}\rangle.
\end{equation}
Indeed, $L_0$ lies in $A_0\otimes B_0$, the $R_s$ lie in distinct
rows of $E\otimes B_0$, and the $S_t$ lie in distinct columns of
$A_0\otimes F$. In particular $\dim L=5+(p-3)+(q-3)=r$.

\begin{lemma}\label{lemma:extended-tangencies}
The intersection of $L$ with the affine rank-one matrix tangent space
at each $Q_i$, $R_s$, and $S_t$ is precisely the scalar line of that
matrix.
\end{lemma}

\begin{proof}
Use the four blocks given by
$(A_0\oplus E)\otimes(B_0\oplus F)$. The space $L$ has no
$E\otimes F$ block.
At a core matrix $Q_i=a_i\otimes b_i$, every extra row of a tangent
matrix is proportional to $b_i$. In $L$, its row indexed by $u_s$
is proportional to $v(s+3)$. These lines differ by
Lemma~\ref{lemma:avoidance}, so all coefficients of the $R_s$ vanish.
In the same way, the extra columns of a tangent matrix are proportional
to $a_i$, whereas those of $L$ are proportional to the corresponding
$v(t+3)$. All coefficients of the $S_t$ vanish. The remaining
intersection in the core block is $\C Q_i$ by
Lemma~\ref{lem:tangent-intersections}.

At $R_s=u_s\otimes v(s+3)$, write a tangent matrix as
\[
 x\otimes v(s+3)+u_s\otimes y.
\]
Its $E\otimes F$ block forces the $F$-part of $y$ to be zero if this
matrix lies in $L$. Its possible $F$-columns are supported on $u_s$,
whereas the $F$-columns of $L$ are supported on $A_0$; hence all
$S_t$ coefficients vanish. For an $E$-row indexed by $u_j$ with
$j\ne s$, the tangent row is proportional to $v(s+3)$ while the
$L$-row is proportional to $v(j+3)$. Their distinct lines force every
other $R_j$ coefficient to be zero. The core block is an element of
$L_0\cap(A_0\otimes\C v(s+3))$, which is zero by
\eqref{eq:avoidance}. Only the row indexed by $u_s$ remains, and its
membership in $L$ makes it a scalar multiple of $R_s$.

At $S_t=v(t+3)\otimes w_t$, interchange rows and columns in this
argument. The $E\otimes F$ block vanishes, the other extra columns
are excluded by distinctness of the $v$-lines, and the core block is
in $L_0\cap(\C v(t+3)\otimes B_0)=0$. The intersection is
$\C S_t$. These comparisons also cover the case when both $E$ and
$F$ are nonzero.
\end{proof}

Let $f_1,\ldots,f_r$ be the coordinate basis of $\C^r$. The required
integer tensor in the seed format $p\times q\times r$ is
\begin{equation}\label{eq:uniform-witness}
 T_{r;p,q}=\sum_{i=1}^5Q_i\otimes f_i
       +\sum_{s=1}^{p-3}R_s\otimes f_{5+s}
       +\sum_{t=1}^{q-3}S_t\otimes f_{p+2+t}.
\end{equation}
Replacing its first five terms by
\begin{equation}\label{eq:uniform-second}
 \sum_{i=1}^4 Q_i\otimes(f_i-f_5)-Q_6\otimes f_5
\end{equation}
and leaving the added terms unchanged gives its second decomposition.
All entries and all displayed decomposition factors are integers.

\begin{proposition}\label{prop:witness-immersion}
Let $r\geq5$, $p,q\geq3$, and $p+q=r+1$.
For any $a\geq p$, $b\geq q$, $c\geq r$, pad
$T_{r;p,q}$ with zero coordinates to format $a\times b\times c$.
Its ordinary rank and border rank are both $r$. The two displayed
unordered $r$-term decompositions are distinct, have independent
third-factor vectors, and have injective addition differentials.
Each differential has rank $r(a+b+c-2)$.
\end{proposition}

\begin{proof}
In either decomposition, the $r$ matrix factors form a basis of $L$:
any five core matrices form a basis of $L_0$, and the added matrices
are independent by~\eqref{eq:extended-L}. The $r$ third factors are
also a basis. In the second decomposition this follows from the
invertible change replacing $f_i$ by $f_i-f_5$ for $1\leq i\leq4$,
with $f_5$ and the remaining basis vectors retained. The choice of
$Q_5$ in the first decomposition and $Q_6$ in the second makes the
unordered pure summands distinct between the two decompositions.
Each summand is nonzero; independence of the third factors makes
summands within a decomposition distinct.

Lemma~\ref{lemma:extended-tangencies} and
Lemma~\ref{lemma:tangent-criterion} prove immersion in the seed
format. The third flattening has image $L$ and rank $r$. Flattening
minors vanish on the border-rank-at-most-$(r-1)$ variety, while the
displayed decomposition has $r$ terms. Thus both tensor rank and
border rank equal $r$.

Choose complements for the three inclusions into the padded spaces.
The derivative directions obtained by varying the first factor in
its complement are
\[
 (\C^a/A')\otimes
   \langle b_i\otimes c_i:1\leq i\leq r\rangle,
\]
with an analogous expression for the second factor; we identify the
quotient with the chosen complement. Since the $c_i$ are independent
and every $a_i,b_i$ is nonzero, both opposite-factor pair lists are
independent. Variation of the third factor in its complement gives
$L\otimes(\C^c/\C^r)$, also with independent matrix factors.
These three spaces occupy different complementary tensor blocks and
are disjoint from the original derivative image. Their dimensions are
$r(a-p)$, $r(b-q)$, and $r(c-r)$. Adding them to the seed rank gives
\[
 r(p+q+r-2)+r(a-p)+r(b-q)+r(c-r)=r(a+b+c-2).
\]
This is the dimension of the product of the $r$ pure-tensor cones in
the padded format. Hence its differential is injective. The free
finite permutation quotient is etale, so the unordered differential
is injective as well. Padding preserves the flattening rank and the
displayed tensor equality, proving all remaining assertions.
\end{proof}

For the hypotheses of Theorem~\ref{thm:main}, choose explicitly
\begin{equation}\label{eq:seed-choice}
 p=\max\{3,r+1-b\},\qquad q=r+1-p.
\end{equation}
Then $p\leq a$ because $a\geq3$ and $a+b\geq r+2$;
also $p\leq r-2$ because $r\geq5$ and $b\geq3$.
It follows that $q\geq3$ and $q\leq b$.
Thus~\eqref{eq:uniform-witness} specifies an integer witness for every
format in the main theorem. At $r=5$ the seed is the original
$3\times3\times5$ tensor. Its two differential ranks after padding
are $50$ in format $3\times4\times5$ and $80$ in format
$6\times6\times6$.
